\clearpage
\SourceLeaf{321}
\section*{REMARKS.}
\pdfbookmark[1]{Remarks on the geometrical aspect of ternary forms}{ternary-bemerkungen}
\BookRule{.15\textwidth}
Besides the preceding notes, which in part develop the ideas given in the review of \textsc{Seeber}'s investigations of ternary forms, the manuscript records several crystal measurements made by Gauss himself with an eight-inch \textsc{Reichenbach} theodolite. Each record gives the date of the observation, showing that this investigation belongs to July $1831$.

Of the theory of indefinite ternary quadratic forms, the manuscripts left behind contain only the following theorem, probably written down while the Disqu. Arr. were being prepared:
\begin{quote}
‘All transformations of the ternary form
\[
\left(\begin{smallmatrix}1,&1,&-1\\0,&0,&0\end{smallmatrix}\right)
\]
into itself are given by the formula
\[
\begin{array}{ccc}
\alpha\delta+\beta\gamma&\alpha\beta-\gamma\delta&\alpha\beta+\gamma\delta\\[.5ex]
\alpha\gamma-\beta\delta&\frac12(\alpha\alpha+\delta\delta-\beta\beta-\gamma\gamma)&\frac12(\alpha\alpha+\gamma\gamma-\beta\beta-\delta\delta)\\[.5ex]
\alpha\gamma+\beta\delta&\frac12(\alpha\alpha+\beta\beta-\gamma\gamma-\delta\delta)&\frac12(\alpha\alpha+\beta\beta+\gamma\gamma+\delta\delta)
\end{array}
\]
with $\alpha,\beta,\gamma,\delta$ chosen so that $\alpha\delta-\beta\gamma=1$.’
\end{quote}

All transformations in which the nine coefficients are integers are obtained by taking for $\alpha,\beta,\gamma,\delta$ both all integers satisfying the condition equation, specifically two even and two odd, and all odd multiples of $\surd\frac12$ satisfying the same condition equation $\alpha\delta-\beta\gamma=1$.

\SourceLeaf{322}
On page $309$. Chaux carbonatée équiaxe, inverse, contrastante and mixte are the names used by \textsc{Hauy} (Traité de Minéralogie $1801$, volume II, pp. $132$, $137$).

The table of transformations of the form $\left(\begin{smallmatrix}5,&5,&5\\1,&1,&1\end{smallmatrix}\right)$ gives, in the first vertical column, the coefficients of the substitution; in the second, the resulting new form; in the third, the primitive form corresponding to the latter, if it is not already primitive; and in the fourth, its adjoint.
\begin{flushright}\textsc{Schering.}\end{flushright}
\BookRule{.30\textwidth}
